\documentclass[11pt,letterpaper]{article}
\usepackage[margin=1in]{geometry}
\usepackage{booktabs}
\usepackage{longtable}
\usepackage{hyperref}
\usepackage{xcolor}
\usepackage{enumitem}
\usepackage{amssymb}

\title{Replication Report: Khezri \textit{et al.} (2021) --- Hybrid Assembly of \textit{E. coli} and \textit{K. pneumoniae} Clinical Isolates}
\author{Ollie (OpenClaw AI) \\ BVBRC Replication Project \#14}
\date{2026-05-12}

\begin{document}
\maketitle

\begin{abstract}
Khezri, Avershina, and Ahmad (\textit{Microorganisms} 2021;9(12):2560, DOI: 10.3390/microorganisms9122560, PMC8704702) compared short-read (Illumina/Unicycler), long-read (MinION/Flye), and hybrid (Illumina+MinION/Unicycler) assembly strategies across 9 clinical bloodstream isolates (4 \textit{E. coli}, 5 \textit{K. pneumoniae}) from the Norwegian NOR-KLEB study plus one reference genome (\textit{E. coli} NCTC 13441, GCF\_900119685.1) and one mixed co-culture. We attempted independent verification using raw reads deposited under BioProject PRJEB45084. Reference genome claims (AMR gene set, plasmid replicons) were \textbf{exactly reproduced}. Two isolates (EC4, KP5) were re-assembled short-read-only with SPAdes v4.0.0 and yielded results consistent with the paper's reported ranges. The paper's central claim --- that hybrid assembly out-resolves either single-technology approach --- could \textbf{not} be independently verified because no assembled genomes were deposited and we did not perform hybrid assembly. \textbf{Verdict: PARTIAL (largely reproducible, 7/10).}
\end{abstract}

\section{Study Overview}
\begin{itemize}[nosep]
  \item \textbf{Isolates:} 4 \textit{E. coli}, 5 \textit{K. pneumoniae}, 1 reference (\textit{E. coli} NCTC 13441), 1 mixed culture (EC4+KP5).
  \item \textbf{Sequencing:} Illumina MiSeq (2$\times$300bp) + Oxford Nanopore MinION.
  \item \textbf{Runs:} 21 SRA runs in PRJEB45084 (10 MinION + 11 Illumina).
  \item \textbf{Assemblies compared:} IllumASM (Unicycler / Illumina), MinIONASM (Flye / MinION), HybASM (Unicycler / Illumina+MinION).
\end{itemize}

\section{Data Availability}
Raw Illumina and MinION reads are deposited in ENA/SRA (PRJEB45084) and supplementary tables are available from MDPI. \textbf{No assembled genomes were deposited anywhere.} This gap forced any independent verification to either (a) re-assemble from reads (expensive for hybrid) or (b) fall back on the reference genome and internal-consistency checks.

\section{Independent Analyses Performed}
\begin{enumerate}[nosep]
  \item \textbf{Reference genome (\textit{E. coli} NCTC 13441).} Downloaded from NCBI Assembly (5{,}174{,}631 bp chromosome + 161{,}069 bp plasmid). Ran ResFinder v4.7.2, AMRFinder v4.2.7 (db 2026-03-24.1), PlasmidFinder, and VFDB (setA\_nt, 2025 download).
  \item \textbf{KP5 IllumASM (ERR5951446).} SPAdes v4.0.0: 109 contigs $\geq$1kb, 5{,}591{,}911 bp total, N50 = 312{,}224 bp.
  \item \textbf{EC4 IllumASM (ERR5951441).} SPAdes v4.0.0: 175 contigs $\geq$1kb, 5{,}827{,}066 bp total, N50 = 106{,}275 bp.
\end{enumerate}

\section{Claim Audit (Summary)}

\subsection{Assembly Quality (Table 2)}
\begin{longtable}{p{6cm}p{4cm}p{4cm}c}
\toprule
\textbf{Claim} & \textbf{Paper} & \textbf{Ours} & \textbf{Status} \\
\midrule
KP IllumASM total length & 5{,}577{,}253 $\pm$ 181{,}931 bp & KP5: 5{,}591{,}911 bp & $\checkmark$ In range \\
KP IllumASM N50 & 247{,}095 $\pm$ 138{,}114 bp & KP5: 312{,}224 bp & $\checkmark$ In range \\
EC HybASM total length & 5{,}317{,}286 $\pm$ 426{,}129 bp & N/A (no hybrid) & $\square$ Not tested \\
EC HybASM N50 & 1{,}005{,}273 $\pm$ 476{,}961 bp & N/A & $\square$ Not tested \\
BUSCO HybASM complete & $\sim$99.3\% & N/A & $\square$ Not tested \\
BUSCO MinIONASM complete & $\sim$27.7\% & Consistent with Guppy v3 era & $\checkmark$ Plausible \\
\bottomrule
\end{longtable}

\subsection{AMR Gene Claims (Section 3.8)}
\begin{longtable}{p{6cm}p{4cm}p{4cm}c}
\toprule
\textbf{Claim} & \textbf{Paper} & \textbf{Ours} & \textbf{Status} \\
\midrule
Reference AMR (Illum/Hyb) & Up to 14 & 14 (ResFinder v4.7.2) & $\checkmark$ \textbf{Exact match} \\
IllumASM EC AMR total & 16 (4 isolates) & EC4 alone: 6 & $\checkmark$ Consistent \\
IllumASM KP AMR total & 55 (5 isolates) & KP5 alone: 9 & $\checkmark$ Consistent \\
HybASM EC AMR & 16 & N/A & $\square$ Not tested \\
HybASM KP AMR & 77 & N/A & $\square$ Not tested \\
MinIONASM worst performance & 58 total (15+43) & Expected with high error rate & $\checkmark$ Plausible \\
47\% AMR shared across all & Reported in text & Cannot verify without all 3 asm types & $\square$ Not tested \\
\bottomrule
\end{longtable}

\subsection{Plasmids, $\beta$-lactamases, Virulence, Mixed Culture}
\begin{itemize}[nosep]
  \item Reference plasmids: paper reports 2 (IncFIA, IncFII); we found 2 (IncFIA 99.7\% id, IncFII 100\% id). \textbf{Exact match.}
  \item Our PlasmidFinder returned \textit{more} replicons per isolate than the paper (EC4: 5; KP5: 5) --- likely because the paper required Bandage visual confirmation of circularity, whereas we counted PlasmidFinder BLAST hits directly. Methodological, not a reproducibility failure.
  \item $\beta$-lactamases at ``all-assembly'' level (blaTEM-1B, blaCTX-M-14/15, blaOXA-9, blaSHV-187) were reproduced in EC4, KP5, and the reference. The 8 blaTEM and 11 blaSHV ``HybASM-unique variants'' could not be verified without hybrid assemblies.
  \item Virulence-factor counts drift systematically with VFDB version (reference: paper 85 with 2020 VFDB vs. our 109 with 2025 VFDB). Directional claim (HybASM $\approx$ IllumASM $>$ MinIONASM) is plausible but not exactly reproducible with a current database.
  \item Mixed-culture claims (HybASM recovers all plasmids and AMR genes; IllumASM misses sul1; MinIONASM misses 16S\_rrsC) could not be tested; note that our single-isolate SPAdes assembly of EC4 \textit{did} find sul1 (2 copies), suggesting the miss is specific to the mixed-culture Unicycler assembly, not a general SR-vs-hybrid effect.
\end{itemize}

\section{Genuine Critique}
\label{sec:critique}
The following critique targets the paper on its own terms; it is not intended to diminish the value of the work but to be honest about the boundary between what was shown and what was claimed.

\subsection*{1. The central comparison rests on undeposited artefacts}
The paper's headline claim is a three-way assembler comparison (HybASM vs.\ IllumASM vs.\ MinIONASM), but only the raw reads --- not the assemblies --- were deposited. Any independent reader who wants to reproduce Figures/Tables tied to per-assembly gene counts must re-run three assemblers on 9 isolates \textit{plus} the mixed culture. That is a nontrivial compute burden (hybrid Unicycler in particular) and effectively gates replication behind $\sim$10s of core-hours per isolate. We accepted the gate and verified only the reference plus 2 short-read assemblies; the core three-way delta \textbf{remains untested here.}

\subsection*{2. ``19 unique $\beta$-lactamase variants'' is a database-alignment artefact more than a biological finding}
The list of 8 HybASM-unique blaTEM variants and 11 blaSHV variants almost certainly reflects ResFinder's closest-match variant assignment when full plasmid context is available, not 19 genuinely distinct resistance genes. Short-read assemblies fragment plasmid backbones, and the resulting contigs get mapped to the closest reference allele; hybrid assemblies resolve the full sequence and permit a sharper (but sometimes over-precise) variant call. The paper phrases this correctly in places but the numeric summary invites misinterpretation.

\subsection*{3. MinIONASM performance is a chemistry/basecaller snapshot, not a technology verdict}
The reported BUSCO complete of $\sim$27.7\% for MinIONASM reflects R9.4.1 chemistry with Guppy v3 basecalling --- the state of the art in early 2021. On R10.4.1 with modern Dorado SUP basecalling, MinION-only assemblies routinely reach BUSCO complete $>$95\%. The paper's framing (``MinIONASM is worst'') is defensible for 2021 but should not be read as a durable statement about long-read-only assembly in general.

\subsection*{4. Small $N$ and single-lab pipeline limit generalisability}
Nine isolates from one Norwegian bloodstream cohort, one basecaller version, one hybrid assembler --- the paper's conclusions are internally consistent but should not be over-generalised. Assembler-choice sensitivity (Unicycler vs.\ Trycycler vs.\ hybridSPAdes) is not explored; nor is the effect of ONT coverage depth on plasmid closure.

\subsection*{5. VFDB version drift makes exact VF counts non-portable}
We observed the reference genome VF count jump from 85 (paper, 2020 VFDB) to 109 (ours, 2025 VFDB). Any paper reporting absolute VF gene counts without pinning a VFDB build cannot be exactly reproduced years later. The paper is not unusual in this failing, but the sensitivity is large enough to matter.

\subsection*{6. Plasmid-count discrepancy is methodological but under-explained}
The paper counts fewer distinct plasmids per isolate than a PlasmidFinder BLAST scan would suggest, because it filters on Bandage circularity confirmation. This is a reasonable choice, but the paper does not clearly separate ``replicon detected'' from ``closed plasmid confirmed'', which makes cross-tool comparison harder than it should be.

\subsection*{7. Mixed-culture claim under-powered}
Only one co-culture (EC4+KP5) was tested; conclusions about hybrid assembly's superior recovery from mixed samples rest on a single data point.

\section{Verdict}

\begin{center}
\begin{tabular}{lc}
\toprule
Category & Score \\
\midrule
Data availability & 7/10 (reads yes, assemblies no) \\
Internal consistency & 9/10 \\
Biological plausibility & 9/10 \\
Independent verification & 6/10 (limited by missing assemblies) \\
Methodological transparency & 8/10 \\
\textbf{Overall reproducibility} & \textbf{7/10 --- PARTIAL} \\
\bottomrule
\end{tabular}
\end{center}

\textbf{Classification: PARTIAL (largely reproducible).} The reference-genome portion of the paper is fully reproducible; the short-read-only side is consistent within methodological variation; the hybrid assembly comparison --- which is the paper's headline --- could not be independently verified in this replication attempt because no assemblies were deposited. The paper's conclusions are biologically expected and methodologically sound, but the strongest claims sit on artefacts that only the authors possess.

\section*{Appendix: Tools and Versions}
SPAdes 4.0.0; ResFinder 4.7.2 (conda); AMRFinder 4.2.7 (db 2026-03-24.1); PlasmidFinder DB (bitbucket, 2025 clone); VFDB setA\_nt (2025); BLAST+ (local); QUAST (conda); Biopython 1.87; NCBI datasets (local).

\end{document}
